% id: 148-14
% book: 베이직쎈 확률과 통계 · 07 통계적 추정
% section: 실전 감각 UP
% figure: tikz:fig-148-14
% answer: 
% answer_source: 
% variant_level: 0 (원본 전사)
% 이 파일은 build.mjs 가 items.json 에서 만든다. 고칠 때는 items.json 을 고친다.
\smash{\raisebox{1.5mm}{\dmkichul{베이직쎈 확률과 통계 148쪽 14번 · 서답형 · 서술형}}}
\noindent\begin{minipage}[t]{0.58\linewidth}\vspace{0pt}\raggedright $\mathrm{A}$ 회사의 무선 청소기 한 대를 최대로 충전했을 때 사용 가능 시간은 평균이 45분, 표준편차가 10분인 정규분포를 따른다고 한다. $\mathrm{A}$ 회사의 무선 청소기 중 임의추출한 $n$대의 사용 가능 시간의 평균을 $\overline{X}$분이라 할 때, 위의 표준정규분포표를 이용하여 $\mathrm{P}(42\le\overline{X}\le 48)=0.4514$를 만족시키는 $n$의 값을 구하시오.\end{minipage}\hfill\begin{minipage}[t]{0.4\linewidth}\vspace{0pt}\centering \resizebox{\ifdim\width>\linewidth\linewidth\else\width\fi}{!}{% 148-14(인쇄 148쪽) — 발문 오른쪽 표준정규분포표($z$ / $\mathrm{P}(0\le Z\le z)$).
% 스캔 표라 크롭하지 않고 tabular 로 다시 조판했다(칸 값은 원문 그대로 · 원문에 없는 행은 넣지 않는다).
{\setlength{\tabcolsep}{4pt}\renewcommand{\arraystretch}{1.3}\small%
\begin{tabular}{|c|c|}
\hline
$z$ & $\mathrm{P}(0\le Z\le z)$ \\
\hline
$0.6$ & $0.2257$ \\
\hline
$0.8$ & $0.2881$ \\
\hline
$1.0$ & $0.3413$ \\
\hline
\end{tabular}}
}\end{minipage}\par
